\scanpage{82}
les $E_i$ qui sont à composantes \unsure{isopyptiques|isopentiques|isopyntiques|isopypiques|isotypiques} $\handwritten{\lambda}$ tels que $0\handwritten{<\lambda\leq\unsure{<}}1$ sont exactement les modules de Dieudonné des groupes formels $p$-divisibles sur $k$ (la catégorie de ces modules de Dieudonné \deleted{correspond à}\inserted{est opposée à} la catégorie de ces groupes, mod isogénie). D'autre part, un F-isocristal est d'amplitude isopentique contenue dans $\handwritten{[}0,1\handwritten{]}$ (intervalle fermé) sss c'est le module de Dieudonné d'un groupe de Barsotti-Tate sur $k$ (la catégorie des F-isocristaux envisagée étant \inserted{anti}équivalente à la catégorie des groupes de B.T. sur $k$ mod. isogénie ..): le F-isocristal unité correspond au groupe de B.T. ind-étale $\handwritten{\underline}{Q}_p/\handwritten{\underline}{Z}_p$ sur $k$.

Supposant toujours $k$ parfait, la condition \deleted{de l'existe} que $(E,F_E)$ soit effectif èt qu'il existe un $V_i$ \inserted{effectif} tel que $FV_i=V_iF=p^i\operatorname{id}$ (cf n\textsuperscript{o}2) signifie que $E$ est d'amplitude isopentique contenue dans $\handwritten{[}0,i\handwritten{]}$.

10. \handwritten{\uline}{Cas} $k$ \handwritten{\uline}{fini}. Supposons $k={\unsure{\handwritten{\underline}{F}|\handwritten{\underline}{\mathbb{F}}}}_q^{\deleted{\handwritten{\unsure{*|\overline{\vphantom{F}}}}}}$, où $q=p^{\deleted{a}\handwritten{\unsure{\nu|a}}}$. En vertu de 6.1, \deleted{son} la gerbe $\handwritten{\underline}{G}(\handwritten{\underline}{F}_q)$ est connue, c'est la sous-gerbe évidente de $\handwritten{\underline}{G}(\handwritten{\underline}{F}_p)$, s'insérant dans une suite exacte
\[
0\handwritten{\longrightarrow}\handwritten{\underline}{G}(\handwritten{\underline}{F}_q)\handwritten{\longrightarrow}\handwritten{\underline}{G}(\handwritten{\underline}{F}_p)\handwritten{\longrightarrow}\operatorname{Rep}(\handwritten{\underline}{Z}/\deleted{a}\handwritten{\nu}\handwritten{\underline}{Z})\handwritten{\longrightarrow}0\quad .
\]
Il est clair alors que pour $\handwritten{\underline}{F}_q$ comme pour $\handwritten{\underline}{F}_p$, la réponse à 6.2 est affirmative, donc les considérations du n\textsuperscript{o}7 s'appliquent: on trouve la structure de $\handwritten{\underline}{G}(\handwritten{\underline}{F}_q)$ comme gerbe + ses structures supplémentaires examinées jusqu'à présent.

Le foncteur fibre \inserted{canonique $(E,F_E)\handwritten{\longmapsto}E$} sur $\operatorname{Fcriso}(\handwritten{\underline}{F}_q)$ \deleted{à valeur} sur le corps $K_{\deleted{a}\handwritten{\nu}}=K(\handwritten{\underline}{F}_{q=p^{\deleted{a}\handwritten{\nu}}})$ \margin{\handwritten{\unsure{présenter|préciser} plutôt $f_q$ comme un \uline{autom. du foncteur identique} de $\operatorname{Fcriso}(F_q)$, donne un élément de $G_\nu(\mathbb{Q}_p)$}}
(i.e. section canonique de $\handwritten{\underline}{G}(\handwritten{\underline}{F}_q)$ sur $K_{\deleted{a}\handwritten{\nu}}$) est muni d'une structure supplémentaire, savoir un automorphisme linéaire $F_E{}^a$, de sorte que le groupe \deleted{algébrique} sur \deleted{k} $K_a$ correspondant à cette section de $\handwritten{\underline}{G}(\handwritten{\underline}{F}_q)$ sur $K_a$, savoir le sous-groupe \inserted{$\handwritten{G_\nu=\nu}$}$\deleted{a}.G_{\deleted{K_a}\handwritten{K_\nu}}$ de $G_{K_{\deleted{a}\handwritten{\nu}}}$ ouvert d'indice $\deleted{a}\handwritten{\nu}$ (correspondant au sous-groupe des racines $\deleted{a}\handwritten{\nu}$.èmes de 1 dans le groupe des caractères $\handwritten{\overline}{\handwritten{\underline}{Q}}^{\handwritten{*}}$ de $G_{\handwritten{K_\nu}}$)

% archive sheet number: 81
